% Proposed application text. Insert into Applications after author review.
% Empirical values come from the committed real-data artifacts; visual concepts remain unverified.
\mypara{Visual Concepts Associated with Human Preference}
We apply our fine-grained dataset-diffing score to recorded human preferences in ImageRewardDB and the original four-candidate Pick-a-Pic release. We score all $54{,}030$ supplied concepts using frozen FG-CLIP2 embeddings and the maximum similarity over valid image patches. ImageReward contributes every strict ranked comparison; Pick-a-Pic contributes the selected image against each non-selected alternative. We average comparisons within annotation events, events within normalized prompts, and prompts equally. Concepts and their directions are selected on discovery data and frozen before held-out evaluation; prompt and image dependencies are kept within splits. The pilot uses a $1{,}024$-patch budget.

The ImageReward pilot retains $487$ valid prompts, including $75$ held-out prompts. Because original Pick-a-Pic image URLs are unavailable, we recover verified image identities from a pinned archive published by the original author. It covers $573$ complete original events across $330$ prompts, requiring $2{,}053$ images; $33$ prompts in $31$ dependency components form the held-out cohort. Original candidate sets and human selections are preserved. This availability-selected January 2023 cohort need not represent the full Pick-a-Pic release.

ImageReward retains the discovery effect direction for $19/20$ selected concepts, whereas Pick-a-Pic retains $10/20$. Pick-a-Pic's preferred-side discovery terms include \emph{lip}, \emph{nose}, and \emph{mouth}, but all ten preferred-side terms reverse effect direction on held-out prompts. Its rejected-side terms include \emph{pavement}, \emph{driveway}, and \emph{carpet}; these retain negative effects, with individual intervals including zero. Direction retention is descriptive and does not establish verified concept presence or simultaneous statistical significance.

As an additional diagnostic, separate from the paper's MLLM discrimination metric, we average discovery-direction-signed concept scores without fitting weights. Held-out paired concordance is $55.1\%$ for ImageReward (95\% component-bootstrap interval $[49.9,59.9]\%$) and $49.2\%$ for Pick-a-Pic ($[37.5,61.1]\%$). Thus the Pick-a-Pic-discovered score does not establish above-chance ordering in this cohort. For secondary transfer, the frozen ImageReward all-pairs list reaches $55.8\%$ ($[43.5,67.7]\%$) on Pick-a-Pic; the ImageReward best-versus-rest list reaches $64.6\%$ ($[53.3,75.8]\%$). We retain both constructions without selecting a protocol on recipient test performance. Across the full bank, held-out signed effect rank correlation between datasets is $-0.044$ for the primary constructions and $-0.035$ using ImageReward best-versus-rest. These descriptive correlations do not test between-dataset differences.

The stronger transfer varies across generator-specific comparison subsets. For ImageReward best-versus-rest concepts evaluated on Pick-a-Pic same-generator pairs, concordance is $47.1\%$ for Stable Diffusion 2.1 ($17$ prompts) and $76.0\%$ for ProtoGen ($16$ prompts). These conditional observation sets differ and do not identify a causal generator effect. Repeated-annotator clustering leaves only nine Pick-a-Pic test components, limiting uncertainty estimation. All recovered images have $1{,}024$ valid scored patches. We report $2{,}000$-resample component-bootstrap intervals without simultaneous guarantees, retain failed replications, and provide supportive examples, counterexamples and random audit pairs. Similarity differences are associations with recorded choices, not calibrated concept prevalence or explanations of why annotators chose an image; visual concept presence still requires human review.
